%@ pages 87-93
%@ status translated
% Original section 9, "Première définition de champ", pp. 87--93.
% The authors abbreviate "n-précatégorie de Segal" as "n-précat de Segal";
% "précat" is not a slip for "préchamp".  So "n-précat" is translated
% "Segal n-precategory" and "n-préchamp" is translated "Segal n-prestack".
% In the first of the two arrays of the proof of Lemma \ref{stuff} the source
% prints $A{p/}$ and $B{p/}$ where the rest of the proof uses $A_{p/}$ and
% $B_{p/}$; the subscript is restored below.
% In the proof of Lemma \ref{grofibchampimplfib} the source prints "Comme $A$
% est $\Gg$-fibrant" at the step where $A'$ is the $\Gg$-fibrant object; the
% print is kept as it stands.
\clearpage
\section{A first definition of stack}
\label{notionchamppage}

We consider again the closed model category structure of
Theorem~\ref{cmf}, relative to a site $\Xx$ with
Grothendieck topology $\Gg$.

Let $A$ be a Segal $n$-prestack.
Let $A\rightarrow A'$ be a trivial cofibration into a $\Gg$-fibrant Segal
$n$-prestack.  We shall say that $A$ is a {\em Segal $n$-stack} (or an {\em
$\infty$-stack}, or indeed just a {\em stack}) if: \newline
---$A(X)$ is a Segal $n$-category for every $X\in \Xx$; and
\newline
---the morphism $A(X)\rightarrow A'(X)$ is an equivalence of
Segal $n$-categories for every $X\in \Xx$.

One sees easily (see the argument of Corollary~\ref{unique} below) that
this condition is independent of the choice of $A\rightarrow A'$.

For $n=0$ the simplicial presheaves which are Segal $0$-stacks are
pre\-cise\-ly those
which satisfy
Thomason's condition of {\em cohomological descent} (see
\cite{TravauxThomason}), or else those which are {\em
flasque
with respect to every object $X$ of the site $\Xx$} in Jardine's sense \cite{Jardine}
(these are also the {\em homotopy sheaves} of
\cite{kobe}
or the {\em flexible sheaves} of \cite{flexible}). Thomason's terminology
``cohomological descent''
is inspired by the cohomological descent of Deligne--Saint-Donat \cite{SGA4b}
and \cite{HodgeIII}.

It is interesting to observe that the following lemmas, analogues of the
classical results on sheaves, are entirely formal consequences of
Quillen's theory \cite{Quillen}
(moreover one can run through the theory of sheaves of sets on a site, as an
essentially easy application of these standard arguments in cmc theory, see
the work of Rezk \cite{Rezk}).

\begin{lemma}
\label{morphentrechamps}
If $f:A\rightarrow B$ is a morphism of Segal $n$-stacks which is a
$\Gg$-weak equivalence, then $f$ is also an object-by-object weak
equivalence (i.e.\ for the coarse topology).
\end{lemma}
{\em Proof:}
By the definition of Segal $n$-stacks we may assume that $A$ and $B$ are
$\Gg$-fibrant. We may factorise the given morphism as
$$
A\rightarrow B' \rightarrow B
$$
where the first morphism is a $\Gg$-trivial cofibration and the second
morphism is a $\Gg$-trivial fibration. A $\Gg$-trivial fibration
possesses the lifting property with respect to all
cofibrations; in
particular it is also trivial for the coarse topology. On the other hand,
since $A$ is $\Gg$-fibrant, there exists a retraction $B'\rightarrow A$
(induced by the identity on $A$). In particular the morphism $A\rightarrow B'$
admits a left inverse. We have thus shown that every morphism
$f:A\rightarrow B$
between Segal $n$-stacks (for $\Gg$) admits a left inverse $g$, in the
homotopy category of Segal $n$-prestacks for the equivalence of the
coarse topology. This applies in
particular to that inverse $g:B\rightarrow A$ (this is where
one sees why we assumed that $B$ is a
stack). We thus obtain a left inverse $h$ of $g$; now $f$ is a right
inverse of $g$, hence $g$ is invertible, which implies that $f$ is
invertible. This allows us to conclude since, according to Quillen
\cite{Quillen}, the only morphisms
whose class is invertible in the homotopy category
are the
equivalences (here, for the coarse topology).
\eop

\begin{lemma}
\label{grofibchampimplfib}
If $A$ is a Segal $n$-prestack then $A$ is $\Gg$-fibrant if and
only
if $A$ is a stack, and is fibrant for the coarse topology.
\end{lemma}
{\em Proof:}
If $A$ is $\Gg$-fibrant then $A$ is
fibrant for the coarse topology, and by definition it is a stack.

Let us therefore suppose that $A$ is fibrant for the coarse topology and
that it is a stack.
Let $A\hookrightarrow A'$ be a $\Gg$-fibrant replacement.
By definition it is a weak equivalence for the coarse topology.
Since $A$ is fibrant for the coarse topology, there exists a retraction
$r:A'\rightarrow A$. Let us then show that $A$ possesses the lifting
property with respect to every $\Gg$-trivial cofibration $X\rightarrow Y$.
Since $A$ is $\Gg$-fibrant, every morphism $X\rightarrow A$ extends
to a morphism $Y\rightarrow A'$. Composing with the retraction $r$
one obtains the desired extension $Y\rightarrow A$.
\eop


{\em Remark:} We have a similar (and even prettier) result for the
HBKQ structure (Theorem~\ref{cmfHirschho}): if $A$ is a
Segal $n$-prestack then $A$ is HBKQ-$\Gg$-fibrant if and
only
if the $A(X)$ are fibrant for every $X$, and $A$ is a stack. In particular,
one obtains the same notion of ``stack'' by using
\ref{cmfHirschho} instead of \ref{cmf}.


\begin{lemma}
\label{stuff}
If $f: A\rightarrow B$ is a
$\Gg$-fibration and if $B$ is a Segal $n$-stack, then $A$ is a
Segal $n$-stack.

If $f: A\rightarrow B$ is a morphism between Segal $n$-stacks
and if $f$ is fibrant for the coarse topology, then
$f$ is $\Gg$-fibrant.
\end{lemma}
{\em Proof:}
For the first part, we begin with the pointwise case, for which we shall
only give a sketch of a proof:
if $f:A\rightarrow B$ is a fibration of Segal $n$-precategories, and if $B$ is
a Segal $n$-category, then $A$ is a Segal $n$-category.
We first remark that the morphisms
$A_{p/}(x_0,\ldots , x_p)\rightarrow B_{p/}(fx_0,\ldots ,fx_p)$ are
fibrations of Segal $n-1$-precategories, so by induction on $n$ we may
assume that the $A_{p/}$ are Segal $n-1$-categories. Now
we consider the diagram $$
\begin{array}{ccc}
A_{p/}&\rightarrow &A_{1/}\times _{A_0}\ldots , \times _{A_0}A_{1/}\\
\downarrow && \downarrow \\
B_{p/}&\rightarrow &B_{1/}\times _{B_0}\ldots , \times _{B_0}B_{1/}.
\end{array}
$$
Let $E\rightarrow A_{1/}\times _{A_0}\ldots , \times _{A_0}A_{1/}$
be a morphism of Segal $n-1$-precategories,
which corresponds to a morphism $U^p(E)\rightarrow A$
(where $U^p(\, )$ is the functor left adjoint to
$A\mapsto A_{1/}\times _{A_0}\ldots , \times
_{A_0}A_{1/}$, which therefore associates to every Segal $n-1$-precategory
a Segal $n$-precategory). We have an inclusion $U^p(E)\rightarrow [p]E$, where $[p]$ is
left adjoint to $A\mapsto A_{p/}$.
This inclusion is a trivial cofibration
of Segal $n$-precategories. Now the fact that $B$ is a Segal
$n$-category implies that there exists an extension of $U^p(E)\rightarrow B$
to a morphism $([p]E)'\rightarrow B$, where $([p]E)'$ is
object-by-object equivalent to $[p]E$ over $\Delta$. The fact that $A\rightarrow B$
is fibrant (combined with the fact that $U^p(E)\rightarrow ([p]E)'$ is a
trivial cofibration) implies that there is likewise an extension
$([p]E)'\rightarrow A$. This implies that the morphism $E\rightarrow
A_{1/}\times _{A_0}\ldots , \times _{A_0}A_{1/}$ lifts, up to homotopy,
to a morphism $E\rightarrow A_{p/}$. One could also give a
relative version of this discussion: given $F\rightarrow A_{p/}$,
$F\hookrightarrow E$ and $E\rightarrow A_{1/}\times _{A_0}\ldots , \times
_{A_0}A_{1/}$, there exists, up to homotopy, an extension to $E\rightarrow
A_{p/}$. One obtains that the Segal map is an equivalence, hence $A$
is a Segal $n$-category.


We now treat the first part for $f:A\rightarrow B$ over
$\Xx$. Let $B\hookrightarrow B'$ be a $\Gg$-fibrant replacement
(which is an object-by-object equivalence), and let
$$
A\hookrightarrow A'\rightarrow B'
$$
be a factorisation in which the second arrow is a fibration for the
coarse
topology and the first is an object-by-object equivalence.
We shall show that the second arrow is a $\Gg$-fibration. Let
$E\rightarrow
E'$ be a $\Gg$-trivial cofibration with a diagram
$$
\begin{array}{ccc}
E&\rightarrow & A'\\
\downarrow && \downarrow \\
E'&\rightarrow & B'.
\end{array}
$$
We may assume that $E\rightarrow A'$ and $E'\rightarrow B'$ are fibrations
for the coarse topology. With this reduction, we have, by the property,
cf.\ \ref{proper}, that
$E\times _{A'}A\rightarrow E$ and $E'\times _{B'}B\rightarrow E'$ are
object-by-object equivalences (we use the fact that the values $B(X)$ and
$B'(X)$ are Segal $n$-categories; likewise for $A(X)$ and $A'(X)$,
for which we use the pointwise version mentioned at the beginning).
On the other hand,
we have a diagram
$$
\begin{array}{ccc}
E\times _{A'}A&\rightarrow & A\\
\downarrow && \downarrow \\
E'\times _{B'}B&\rightarrow & B,
\end{array}
$$
in which the left-hand vertical arrow is a $\Gg$-trivial cofibration.
Since by hypothesis the right-hand arrow is a $\Gg$-fibration,
we obtain a lift $E'\times _{B'}B\rightarrow A$. Now
$$
(E'\times _{B'}B)\cup ^{E\times _{A'}A}E\rightarrow E'
$$
is an equivalence for the coarse topology.
We factorise it as
$$
(E'\times _{B'}B)\cup ^{E\times _{A'}A}E\rightarrow E''\rightarrow E'
$$
where the first arrow is a trivial cofibration for the coarse
topology and the second is a trivial fibration for the
coarse topology. We have a morphism coming from our lift
$$
(E'\times _{B'}B) \cup ^{E\times _{A'}A}E \rightarrow A'.
$$
Since $A'\rightarrow B'$ is a fibration for the
coarse topology, there exists an extension $E''\rightarrow A'$.
Now, since $E''\rightarrow E'$ is a trivial fibration for
the coarse topology, there exists a section $E'\rightarrow E''$
which agrees on $E$ with the given morphism $E\rightarrow E''$.
We thus obtain, by composition, a morphism $E'\rightarrow A'$ which provides the
desired lift.



For the second part of the lemma,
we have a morphism between Segal $n$-stacks $f: A\rightarrow B$ which is
fibrant for the coarse topology. We factorise it as
$$
A\stackrel{g}{\rightarrow} A'\stackrel{h}{\rightarrow} B
$$
with $g$ a $\Gg$-equivalence and $h$ a $\Gg$-fibration. By the
first part,
$A'$ is a stack. By the preceding lemma, $g$ is an equivalence for the
coarse topology. If $E\rightarrow E'$ is a $\Gg$-trivial cofibration
with a diagram
$$
\begin{array}{ccc}
E&\rightarrow & A\\
\downarrow && \downarrow \\
E'&\rightarrow & B
\end{array}
$$
then there exists an intermediate lift $E'\rightarrow A'$. As
above, we may assume that the morphism $E'\rightarrow A'$ is a
fibration
for the coarse topology. In this case $E'\times _{A'} A\rightarrow E'$
is an object-by-object equivalence (here we again use the fact that the
$A(X)$ are Segal $n$-categories, cf.\ the beginning of the
proof),
so it is a trivial cofibration for the coarse topology.
We have the inclusion
$E\subset E'\times _{A'}A$. We then apply the lifting
property of $f$ (which is a
fibration for the coarse topology) to the diagram
$$
\begin{array}{ccc}
E'\times _{A'}A&\rightarrow & A\\
\downarrow && \downarrow \\
E'&\rightarrow & B,
\end{array}
$$
to obtain a lift $E'\rightarrow A$ which answers the question.
\eop

\subnumero{The associated stack}


If $A$ is a Segal $n$-prestack on $\Xx$ we shall call an {\em associated
stack of $A$} any weak
equivalence
$A\rightarrow A'$ where $A'$ is a Segal $n$-stack. We note that an associated
stack always exists: it suffices to take for $A'$ a fibrant replacement of
$A$.

\begin{corollary}
\label{unique}
Let $A$ be an $n$-prestack and let $A\rightarrow A'$ and $A\rightarrow A''$
be two stacks associated to $A$. Suppose that $A''$ is fibrant for the
coarse topology. Then there exists an equivalence (for the coarse topology)
$A'\rightarrow A''$ making the triangle commute up to homotopy.
If moreover $A\rightarrow A'$ is a cofibration, one may suppose
the triangle
strictly commutative.
\end{corollary}
{\em Proof:}
If $A\rightarrow A'$ is a cofibration, then it is a trivial cofibration and
the fact that $A''$ is fibrant implies the existence of the desired
extension $A'\rightarrow
A''$.  If $A\rightarrow A'$ is not a cofibration, we factorise
$$
A\rightarrow B \rightarrow A'
$$
with a trivial cofibration on the left and a trivial fibration
on the right. There therefore exists a section $A'\rightarrow B$ as well as an extension
of the morphism $A\rightarrow A''$ to $B\rightarrow A''$. This gives a morphism
$A'\rightarrow A''$ and one checks (following the definitions of
\cite{Quillen})
that this morphism makes the triangle commute up to homotopy.
\eop

{\em Improvement:} \,\, With the internal $\underline{Hom}$ of
fibrant Segal
$n$-prestacks which will be defined in \S 11 below, we shall be able to formulate a
better statement of uniqueness of the ``associated stack'' via the following
universal property. If $A\rightarrow A'$ is
a stack associated to the prestack $A$, then for every Segal $n$-prestack $B$
fibrant for the coarse topology and which is a stack, the induced
morphism
$$
\underline{Hom}(A',B)\rightarrow \underline{Hom}(A,B)
$$
is an equivalence of (fibrant) Segal $n$-stacks.
See \S 13.



\subnumero{Fibre products}

Let $A,B,C$ be Segal $n$-precategories. If $B\rightarrow A$
and $C\rightarrow A$ are morphisms, we have at our disposal a {\em homotopy
fibre product} $B\times ^h_AC$, essentially well defined: one may define it as
$B'\times _{A'}C'$ where
$$
\begin{array}{ccccc}
B&\rightarrow & A & \leftarrow & C \\
\downarrow && \downarrow && \downarrow \\
B'&\rightarrow & A' & \leftarrow & C'
\end{array}
$$
is a diagram in which $A',B'$ and $C'$ are Segal
$n$-categories with
$B'\rightarrow A'$ fibrant, and in which the vertical arrows are
weak
equivalences. It is shown in \cite{limits} Lemma 6.1.2, that this product is a
limit in the Segal $(n+1)$-category $nSeCAT$ (\cite{limits} treats the case of
$n$-categories but the case of Segal $n$-categories is identical).

The replacement of $A,B,C$ by
$A',B',C'$ may be chosen functorially. Hence, if $B\rightarrow
A\leftarrow C$ is a diagram of Segal $n$-prestacks on a category
$\Xx$, one may consider their {\em object-by-object homotopy fibre
product}
$B\times ^h_AC$. It is an essentially well-defined $n$-prestack on
$\Xx$ and (with the notations of internal structure which will be introduced in
\S 11
below) it is the limit in $nSeCHAMP (\Xx ^{\rm gro})$ of the
given diagram.

The following lemma is the analogue of the statement in sheaf theory according to which
the kernel of a morphism of sheaves of groups (or the kernel of a double
arrow, in the case of sheaves of sets) is again a sheaf.

\begin{lemma}
\label{fiprodchamp}
Let $\Xx , \Gg$ be a site and $B\rightarrow A\leftarrow C$ a
diagram of Segal $n$-stacks on $\Xx$. Then
the object-by-object homotopy fibre product $B\times ^h_AC$ is a Segal
$n$-stack (and it is also the limit of the diagram in $nSeCHAMP(\Xx ^{\Gg})$,
see the
notation in \S 11).
\end{lemma}
{\em Proof:}
One may choose a replacement of the diagram
$$
\begin{array}{ccccc}
B&\rightarrow & A & \leftarrow & C \\
\downarrow && \downarrow && \downarrow \\
B'&\rightarrow & A' & \leftarrow & C'
\end{array}
$$
with $A', B', C'$ $\Gg$-fibrant, and in which the vertical arrows are
weak equivalences for $\Gg$, and $B'\rightarrow A'$ is a
$\Gg$-fibration. By Lemma~\ref{morphentrechamps} the vertical
arrows are object-by-object weak equivalences, and we therefore have
an adequate replacement for the
construction of the homotopy fibre product
$$
B\times ^h_AC = B'\times _{A'}C'
$$
which is $\Gg$-fibrant, in particular it is a Segal $n$-stack. Any
other construction of $B\times ^h_AC$ will give an $n$-prestack
object-by-object equivalent to the preceding one; and by definition, it will
always be a
Segal $n$-stack.
\eop

The following corollary answers the question originally posed by C.~Rezk
in \cite{Rezk} (where he gives the proof for
$0$-prestacks): passing to the ``associated stack'' is compatible
with homotopy fibre products.

\begin{corollary}
\label{rezkQ}
Let $B\rightarrow A \leftarrow C$ be a diagram of Segal $n$-prestacks.
Suppose that
the values $A(X)$, $B(X)$ and $C(X)$ are Segal $n$-categories.
Let $B'\rightarrow A'\leftarrow C'$ be the diagram obtained by passing
to the associated $\Gg$-stacks. Then
$$
B\times ^h_AC\rightarrow B'\times ^h_{A'}C'
$$
is an associated $\Gg$-stack; here the homotopy fibre products on the two
sides are taken object-by-object over $\Xx$.
\end{corollary}
{\em Proof:}
The arrow in question is a $\Gg$-weak equivalence by Corollary
\ref{cordereedy}; and by the preceding lemma,
the target is a $\Gg$-stack. It
follows by definition that the arrow is an ``associated stack''.
\eop
